\chapter{The HTTP Protocol (Implemented Subset)}
\label{chap:http-protocol}

\section{Scope of This Implementation}

Full HTTP/1.1 (RFC~9110--9112) is a large specification covering persistent connections, chunked transfer encoding, content negotiation, caching, range requests, and many other features. This project implements a deliberately narrow, but protocol-correct, subset:

\begin{itemize}
    \item Only the \textbf{GET} method is supported; other methods receive a \texttt{405 Method Not Allowed} response.
    \item Every response includes \texttt{Connection: close}, meaning each TCP connection handles exactly one request/response pair --- there is no persistent-connection (keep-alive) support.
    \item Content length is always communicated via an explicit \texttt{Content-Length} header; chunked transfer encoding is not implemented.
    \item A minimal set of status codes is used: \texttt{200 OK}, \texttt{400 Bad Request}, \texttt{404 Not Found}, and \texttt{405 Method Not Allowed}.
\end{itemize}

This scope was chosen to cover the concepts that most clearly extend beyond what Gopher already teaches (Chapter~\ref{chap:gopher-protocol}), without requiring an implementation of the full specification.

\section{Request Format}

\begin{lstlisting}[style=bashstyle, numbers=none]
GET /index.html HTTP/1.1
Host: localhost
Connection: close

\end{lstlisting}

An HTTP request consists of three parts, each terminated by \texttt{\textbackslash r\textbackslash n}:

\begin{enumerate}
    \item \textbf{Request line}: \texttt{METHOD SP path SP HTTP-version}, e.g.\ \texttt{GET /index.html HTTP/1.1}.
    \item \textbf{Headers}: zero or more \texttt{Name: value} lines.
    \item \textbf{A blank line} (i.e.\ two consecutive \texttt{\textbackslash r\textbackslash n} sequences) marking the end of the headers and the start of any body.
\end{enumerate}

For a \texttt{GET} request, there is normally no body, so the blank line effectively marks the end of the entire request.

\section{Response Format}

\begin{lstlisting}[style=bashstyle, numbers=none]
HTTP/1.1 200 OK
Content-Type: text/html
Content-Length: 137
Connection: close

<html>...</html>
\end{lstlisting}

The response mirrors the request's structure:

\begin{enumerate}
    \item \textbf{Status line}: \texttt{HTTP-version SP status-code SP reason-phrase}, e.g.\ \texttt{HTTP/1.1 200 OK}.
    \item \textbf{Headers}, most importantly \texttt{Content-Type} (telling the client how to interpret the body) and \texttt{Content-Length} (telling the client exactly how many bytes of body follow).
    \item \textbf{A blank line}.
    \item \textbf{The body} --- exactly \texttt{Content-Length} bytes.
\end{enumerate}

\section{The Central Technical Challenge: Framing}
\label{sec:http-framing}

The hardest part of implementing HTTP correctly, and the concept this project is specifically designed to surface, is: \emph{how does a reader know when the message is complete?}

Unlike Gopher, where the answer is simply ``when the connection closes'' (Chapter~\ref{chap:gopher-protocol}), HTTP as implemented here requires two distinct steps:

\begin{enumerate}
    \item \textbf{Find the end of the headers.} This is always marked by the four-byte sequence \texttt{\textbackslash r\textbackslash n\textbackslash r\textbackslash n}. Everything before it is the status/request line plus headers; everything after it is the body.
    \item \textbf{Determine the body length.} Parse the \texttt{Content-Length} header (if present) from the headers section, and read exactly that many bytes following the blank line.
\end{enumerate}

This project's HTTP client (Chapter~\ref{chap:http-client}) implements both steps explicitly, since --- unlike the server, which \emph{generates} a well-formed response and therefore always knows its own body length --- the client must \emph{parse} an arbitrary incoming response without knowing its length in advance.

\section{Comparison with Gopher}

\begin{center}
\begin{tabularx}{\textwidth}{l X X}
\toprule
\textbf{Aspect} & \textbf{HTTP (this implementation)} & \textbf{Gopher} \\
\midrule
Request structure & Request line + headers + blank line & Single selector line \\
Response structure & Status line + headers + blank line + body & Raw content \\
End-of-message signal & \texttt{Content-Length} header & Connection closure \\
Status reporting & Numeric codes (200, 400, 404, 405) & None (informal type-\texttt{3} error lines) \\
Methods & GET (others rejected with 405) & None --- always ``fetch'' \\
Content typing & \texttt{Content-Type} header, guessed from file extension & None --- client must infer from context \\
\bottomrule
\end{tabularx}
\end{center}

This table extends the one presented in Chapter~\ref{chap:gopher-protocol}, and is the clearest single summary of what implementing HTTP \emph{after} Gopher actually adds, conceptually.
